
\documentclass[11pt,a4paper]{article}

\usepackage[T1]{fontenc}
\usepackage[utf8]{inputenc}
\usepackage[ngerman]{babel}
\usepackage{lmodern}
\usepackage{microtype}
\usepackage{amsmath,amssymb,amsthm}
\usepackage{mathtools}
\usepackage{geometry}
\usepackage{booktabs}
\usepackage{array}
\usepackage{longtable}
\usepackage{graphicx}
\DeclareGraphicsExtensions{.pdf,.png,.jpg}
\graphicspath{{./}{fig/}{figs/}{images/}{img/}}

\newcommand{\includegraphicssafe}[2][]{%
  \IfFileExists{#2}{\includegraphics[#1]{#2}}{%
    \fbox{\parbox{0.9\linewidth}{\centering Abbildung nicht gefunden:\\\texttt{#2}}}%
  }%
}

\geometry{margin=2.7cm}

\theoremstyle{plain}
\newtheorem{satz}{Satz}
\newtheorem{lemma}{Lemma}
\newtheorem{korollar}{Korollar}

\theoremstyle{definition}
\newtheorem{definition}{Definition}
\newtheorem{beispiel}{Beispiel}

\theoremstyle{remark}
\newtheorem{bemerkung}{Bemerkung}

\title{Das offene arithmetische Periodensystem\\\large Eine erste strukturierte Rohfassung}
\author{Thomas Hoffbauer}
\date{\today}

\begin{document}
\maketitle

\begin{abstract}
Wir schlagen eine erste formale Fassung eines \emph{offenen arithmetischen Periodensystems} der natürlichen Zahlen vor. 
Die Konstruktion basiert auf den Basisschalen der Primzahlen \(2,3,5\), dem offenen Restgenerator \(R_{235}(N)\), einer internen Vierergruppenstruktur \(\{e,a,b,c\}\), sowie vier modulierenden Strukturgrößen, die wir heuristisch mit den Motiven von von Klitzing, Aharonov--Bohm, Morley--Walter und Kepler verbinden. 
Im Unterschied zum chemischen Periodensystem ist das hier beschriebene Schema nicht abgeschlossen, sondern wächst mit dem Auftreten neuer Restgeneratoren offen nach oben weiter. 
Eine verfeinerte Tafel für \(N\le 300\) illustriert die Trennung in Generatorfamilien, Nachläufer, Plateauzustände und Mischstrukturen.
\end{abstract}

\section{Einleitung}

Die natürlichen Zahlen werden in der klassischen Zahlentheorie meist als linear geordnete Folge betrachtet. 
Im vorliegenden Zugang schlagen wir vor, sie als Zustände eines \emph{offenen arithmetischen Periodensystems} zu lesen. 
Die Rede von einem Periodensystem ist dabei nicht im Sinn einer endlichen Tabelle zu verstehen, sondern als Hinweis auf eine schalenartige, gruppierte und nach oben offene Ordnungsstruktur.

Die leitende Intuition lautet:
\begin{quote}
Jede natürliche Zahl \(N\) besitzt eine mehrschichtige arithmetische Signatur, die sich aus Basisschalen, Restgeneratoren, internen Klassen und modulierenden Strukturgrößen zusammensetzt.
\end{quote}

Der hier entwickelte Ansatz ist eine erste Rohfassung. 
Er erhebt keinen Anspruch auf eine abgeschlossene axiomatische Theorie, sondern will eine belastbare Klassifikationssprache bereitstellen, in der Zahlen nicht nur faktoriell, sondern zugleich schalenartig, gruppiert und generatoroffen gelesen werden können.

\section{Basisschalen und Restgeneratoren}

\begin{definition}[Basisschalenkoordinaten]
Für \(N\in\mathbb{N}\) definieren wir die Basisschalenkoordinaten durch die p-adischen Bewertungen
\[
\nu_2(N),\qquad \nu_3(N),\qquad \nu_5(N).
\]
Sie beschreiben die Einbettung von \(N\) in die dyadischen, triadischen und pentadischen Basisschalen.
\end{definition}

Damit erhält jede Zahl zunächst eine dreifache Schalenadresse in der durch \(2,3,5\) erzeugten Basisschalenwelt.

\begin{definition}[2-3-5-Rest]
Für \(N\in\mathbb{N}\) definieren wir den 2-3-5-Rest durch
\[
R_{235}(N):=\frac{N}{2^{\nu_2(N)}3^{\nu_3(N)}5^{\nu_5(N)}}.
\]
\end{definition}

\begin{bemerkung}
Die Bedingung \(R_{235}(N)=1\) charakterisiert genau jene Zahlen, die vollständig in der Basisschalenwelt der Primzahlen \(2,3,5\) liegen. 
Für \(R_{235}(N)>1\) tritt offene Generatorstruktur auf.
\end{bemerkung}

\begin{lemma}
Für jede natürliche Zahl \(N\) gilt die eindeutige Zerlegung
\[
N=2^{\nu_2(N)}3^{\nu_3(N)}5^{\nu_5(N)}R_{235}(N).
\]
\end{lemma}

\begin{proof}
Dies folgt unmittelbar aus der Definition der p-adischen Bewertungen.
\end{proof}

\section{Interne Vierergruppenstruktur}

\begin{definition}[EABC-Klasse]
Sei \(\operatorname{strip}_{2,3}(N)\) die Zahl, die aus \(N\) durch Herausziehen aller Faktoren \(2\) und \(3\) entsteht. 
Die EABC-Klasse \(u(N)\) wird durch die modulo-\(12\)-Reduktion dieses Restes bestimmt:
\[
1\mapsto e,\qquad 5\mapsto a,\qquad 7\mapsto b,\qquad 11\mapsto c.
\]
Damit ist
\[
u(N)\in\{e,a,b,c\}\cong V_4,
\]
wobei \(V_4\) die Kleinsche Vierergruppe bezeichnet.
\end{definition}

\begin{bemerkung}
Die EABC-Klasse ersetzt keine Primfaktorinformation. 
Sie ist als interne Familienvariable zu lesen und ergänzt die Basisschalenkoordinaten um eine diskrete Symmetrieschicht.
\end{bemerkung}

\section{Vier Modifikationsmodule}

Zur Basisschalen- und Reststruktur treten vier modulierende Operatoren hinzu.

\subsection*{(i) von Klitzing}
Das von-Klitzing-Modul steht für robuste Plateau-Bildung. 
Es modelliert die Tendenz, dass arithmetische Zustände sich nicht beliebig fein ändern, sondern diskrete Stufen oder Niveaus ausbilden.

\subsection*{(ii) Aharonov--Bohm}
Das Aharonov--Bohm-Modul liefert eine arithmetische Phase. 
Es beschreibt, dass Zahlen nicht nur lokal durch Faktorisierung, sondern auch durch zyklische Umlaufstruktur charakterisiert werden können.

\subsection*{(iii) Morley--Walter}
Das Morley--Walter-Modul erfasst geometrische Kopplung zwischen den Schalenexponenten, insbesondere dann, wenn mehrere Basisschalen zugleich aktiv sind.

\subsection*{(iv) Kepler}
Das Kepler-Modul misst Harmonie bzw.\ Einbettung einer Zahl in die globale Schalenordnung, also ihre Stärke der Integration in die Basisschalenwelt.

\section{Erste diskrete Realisierung}

\begin{definition}[Arithmetischer Phasenindex]
Wir definieren einen ersten diskreten Phasenindex durch
\[
\phi_{\mathrm{AB}}^{(0)}(N)
=
\nu_2(N)+2\nu_3(N)+3\nu_5(N)+\chi_4(u(N)),
\]
wobei
\[
\chi_4(e)=0,\qquad
\chi_4(a)=1,\qquad
\chi_4(b)=2,\qquad
\chi_4(c)=3.
\]
\end{definition}

\begin{definition}[Morley--Walter-Kopplung]
Die Morley--Walter-Kopplung sei
\[
M_{\mathrm{MW}}(N)
=
\nu_2(N)\nu_3(N)+\nu_3(N)\nu_5(N)+\nu_2(N)\nu_5(N).
\]
\end{definition}

\begin{definition}[Kepler-Harmonie]
Die Kepler-Harmoniezahl sei
\[
K_{\mathrm{Kep}}(N):=\nu_2(N)+\nu_3(N)+\nu_5(N).
\]
\end{definition}

\begin{definition}[von-Klitzing-Plateau]
Als erste diskrete Plateaugröße setzen wir
\[
Q_{\mathrm{vK}}(N):=M_{\mathrm{MW}}(N)+K_{\mathrm{Kep}}(N).
\]
\end{definition}

\begin{lemma}
Ist \(N\) von der Form \(2^n\), \(3^m\) oder \(5^o\), so gilt
\[
M_{\mathrm{MW}}(N)=0.
\]
\end{lemma}

\begin{proof}
In diesem Fall ist höchstens eine der Größen \(\nu_2(N),\nu_3(N),\nu_5(N)\) positiv. 
Damit verschwinden alle Produkte
\[
\nu_2(N)\nu_3(N),\qquad \nu_3(N)\nu_5(N),\qquad \nu_2(N)\nu_5(N).
\]
\end{proof}

\section{Zustandsvektor und Statusklassen}

\begin{definition}[Arithmetischer Zustandsvektor]
Zu jeder Zahl \(N\in\mathbb{N}\) definieren wir den Zustandsvektor
\[
\mathcal P(N)
=
\bigl(
\nu_2(N),\nu_3(N),\nu_5(N),
R_{235}(N),
u(N),
\phi_{\mathrm{AB}}^{(0)}(N),
M_{\mathrm{MW}}(N),
K_{\mathrm{Kep}}(N),
Q_{\mathrm{vK}}(N)
\bigr).
\]
\end{definition}

\begin{definition}[Statusoperator]
Zusätzlich ordnen wir jeder Zahl einen Status
\[
\Sigma(N)\in
\{
\text{Basis},
\text{a-Schale},
\text{Misch-Basis},
\text{Plateau},
\text{Plateauwechsel},
\text{neuer Generator},
\text{harmonischer Nachläufer},
\text{Mischzahl}
\}
\]
zu.
\end{definition}

\begin{bemerkung}
Die Statusklassen sind phänomenologisch zu lesen:
\begin{itemize}
    \item \textbf{Basis}: reine 2-3-5-Welt,
    \item \textbf{a-Schale}: pentadische Spezialschicht,
    \item \textbf{Misch-Basis}: mehrere Basisschalen zugleich aktiv,
    \item \textbf{Plateau / Plateauwechsel}: robuste Stufenstruktur,
    \item \textbf{neuer Generator}: erstmals auftretender Restprimträger \(>5\),
    \item \textbf{harmonischer Nachläufer}: Verstärkung eines bereits vorhandenen Restgenerators,
    \item \textbf{Mischzahl}: hybride Rest- und Schalenstruktur.
\end{itemize}
\end{bemerkung}

\section{Definition des offenen arithmetischen Periodensystems}

\begin{definition}[Offenes arithmetisches Periodensystem]
Das \emph{offene arithmetische Periodensystem} ist die Klassifikation der natürlichen Zahlen \(N\in\mathbb{N}\) durch den Zustandsvektor \(\mathcal P(N)\), ergänzt um den Statusoperator \(\Sigma(N)\).
\end{definition}

\begin{satz}[Offene Periodizität]
Das offene arithmetische Periodensystem ist periodisch im strukturellen, aber nicht im endlichen Sinn:
\begin{enumerate}
    \item Die Basisschalen erzeugen wiederkehrende Schalenmuster.
    \item Die EABC-Klassen erzeugen interne Familien.
    \item Neue Restgeneratoren erweitern das System nach oben offen.
\end{enumerate}
\end{satz}

\begin{proof}[Begründung]
Punkt (1) folgt aus der wiederholten Aktivierung der Basisschalen durch \(\nu_2,\nu_3,\nu_5\). 
Punkt (2) folgt aus der diskreten Vierergruppenstruktur \(u(N)\in\{e,a,b,c\}\). 
Punkt (3) folgt daraus, dass mit neuen Primträgern \(>5\) stets weitere Restgeneratorfamilien auftreten, sodass keine endliche Schließung des Systems vorliegt.
\end{proof}

\begin{bemerkung}[Mendelejew-Analogie]
Die Analogie zum Periodensystem beruht auf vier Eigenschaften:
\begin{enumerate}
    \item Ordnung bekannter Zustände,
    \item Zerfall in Gruppen und Perioden,
    \item Sichtbarkeit von Übergängen und Lücken,
    \item offene Erweiterbarkeit durch neue Generatoren.
\end{enumerate}
In diesem Sinn kann das hier vorgeschlagene Schema als Mendelejew-System der natürlichen Zahlen gelesen werden.
\end{bemerkung}

\section{Erste empirische Tafel}

Für \(N\le 300\) ergibt sich bereits eine verfeinerte Tafel, in der die offenen Generatoren in Familien nach EABC-Klasse zerfallen. 
Insbesondere treten Zeilen der Form
\[
e:\mathrm{Gen},\qquad a:\mathrm{Gen},\qquad b:\mathrm{Gen},\qquad c:\mathrm{Gen}
\]
auf, während Plateaus, Nachläufer und Mischzahlen als Sekundärzustände sichtbar werden.

\begin{figure}[t]
    \centering
    \includegraphicssafe[width=\textwidth]{offenes_arithmetisches_mendelejew_tafel_bis_300_verfeinert.png}
    \caption{Verfeinerte Mendelejew-Tafel des offenen arithmetischen Periodensystems für \(N\le 300\). Die Spalten kodieren verdichtete Schalenbereiche, die Zeilen EABC-Klasse und Statusfamilie.}
    \label{fig:open_arithmetic_periodic_table_refined}
\end{figure}

\section{Interpretation}

Die verfeinerte Tafel legt nahe, dass die natürlichen Zahlen nicht nur nach Faktorisierung, sondern zugleich nach
\[
\text{Schale},\quad \text{Gruppe},\quad \text{Restgenerator},\quad \text{Plateau},\quad \text{Mischungsstatus}
\]
geordnet gelesen werden können.

Die Generatoren erscheinen dabei als offene Erweiterung des Systems, während die Basisschalen und Vierergruppenklassen die lokale Regelmäßigkeit liefern. 
In dieser Hinsicht verhält sich das vorgeschlagene Schema eher wie eine \emph{arithmetische Chemie} als wie eine bloße alternative Notation der Primfaktorzerlegung.

\section{Ausblick}

Die vorliegende Fassung ist als erste diskrete Rohform zu verstehen. 
Naheliegende nächste Schritte sind:
\begin{itemize}
    \item Einbau einer Triggerfunktion \(\Theta(N)\) oder \(\Pi_{\mathrm{open}}(N)\),
    \item Erweiterung auf größere Bereiche \(N\le 1000\) oder \(N\le 5000\),
    \item Kopplung an BM-Datenquellen, AHPN-Listen oder Generatorhistorien,
    \item Verfeinerung der Plateau-, Phasen- und Geometrieoperatoren.
\end{itemize}

Damit könnte aus dem hier vorgeschlagenen Schema ein systematisches arithmetisches Ordnungsmodell entstehen.

\end{document}
